%% Jouni Ritvanen, jouni.ritvanen@lut.fi
%% This section adds chapter 4
\section{MVA Development Process}

Based on the literature review, a MVA development process has been built on top of the recent research and practical insight. As a base for the process the general model for software architecture design process by \textcite{hofmeisterGeneralModelSoftware2007} is used. The process is divided into three main sections that can be found from figure \ref{fig:Architecture_design_activities_hofmeister}. The first main section is Architectural Analysis, the second is Architectural Synthesis, and the third is Architectural Evaluation. However, according to \citeauthor{hofmeisterGeneralModelSoftware2007} these steps are not sequential but happening simultaneously in small leaps. A concept of backlog is utilized. It includes the architectural problems that need to be addressed and can be modified any time during the design process.

Another fundamental insight impacting and guiding the whole process is by \textcite{terragniFutureAIDrivenSoftware2025} who argue that software engineers and AI will work in close symbiosis. As presented in figure \ref{fig:AI_in_development_process_terragni}, there will be a constant communication flowing between the AI and the software engineer. AI will generate and verify the things it has produced. AI will also explain to the software engineer what it has done and why. The engineer will give AI initial and guiding instructions, and enquiries regarding the outputs. The engineer will also improve and verify the outputs given by AI in every step. According to the model, AI is never working fully autonomously without the guidance and verification of the engineer. However, it also seems like engineer is always utilizing the AI in some suitable way. This insight will also guide the MVA design process.

Third thing affecting the process are the characteristics of software start-ups. They have very limited resources, are often young and inexperienced, try to get into new markets using new technologies, and have multiple actors influencing them \parencite{suttonRoleProcessSoftwareJuly-Aug./2000}. Especially the resources affect the MVA design process, since it needs to be relatively light, cheap, and fast. Software start-ups don't often have too much time, money, and workforce available. This is another reason why AI will be largely used in different parts of the process: the productivity of a software engineer can be increased by 35\% using AI tools \parencite{amasantiImpactAIGeneratedSolutions2026}.

\subsection{Architectural Analysis}

The purpose of Architectural Analysis is to figure out the set of ASRs for the system.In the general model provided by \textcite{hofmeisterGeneralModelSoftware2007}, Architectural Analysis has two inputs: architectural concerns and context. Architectural concerns include requirement, but they might also include things like mandatory design decisions such as standards. The context means the environment of the system. That includes "the developmental, operational, political, and other influences of the system." In this process it is assumed, that these two inputs have already been formed.

As a first step, these documents are given to the chosen LLM tool. It is prompted to analyze the context and requirements to find out the ASRs based on the list of indicators by \textcite{chenCharacterizingArchitecturallySignificant2013}. Examples of all prompts will be given in the demonstration in section 5. The LLM is asked to provide a list of ASRs it gathered from the given documents. The list must include the requirements, arguments for why each requirement was recognised to be significant, and some indication of how sure the LLM is of each ASR. The LLM is also asked to provide a list of requirements it is not sure whether or not they are architecturally significant.

A study made by \textcite{tagliaferroLeveragingLLMsAutomate2025} found that LLMs are able to produce UML components from natural language specifications. In this case smaller steps are taken at a time, since that is more in line with the constant communication flow proposed by \textcite{terragniFutureAIDrivenSoftware2025}. All of these outputs are asked, so that the engineer can go through what the AI has produced, and why it came to different conclusions. This gives the engineer the opportunity to correct mistakes made by the LLM in an earlier stage, guide it properly in the next stages, and see if there is more context, input, or information needed. The engineer can do their own analysis, add more ARSs, remove some or make decisions in the case the LLM is not sure whether a requirement is significant or not. Because these phases aren't sequential, both the LLM and the engineer can modify the list of ASRs any time during the design process.

\subsection{Architectural Synthesis}

According to \textcite{hofmeisterGeneralModelSoftware2007}, the purpose of this design phase is to propose solutions for the ASRs specified in the previous phase. Requirement-to-architecture is the part of software development where LLMs are most often utilized \parencite{espositoGenerativeAISoftware2026}.

At this phase, the input for LLM are the ASRs. The LLM is prompted to generate multiple different ideas with the goal of ensuring that all of the ASRs are fulfilled. These ideas can be certain technologies, containers, components, modules, or other solutions to the problems. Each ASR gets its own list of possible solutions. All of the ideas must be justified and presented with pros, cons, and possible trade-offs. \textcite{cervantesLLMsAssistantsSoftware2026} found that LLMs are very useful in this kind of divergent thinking. LLMs are able to generate a massive amount of ideas and solutions with justifications. However, they lack the skills in convergent thinking, e.g. putting the different ideas together.

Therefore the engineer is presented with the list of possible solutions and given the option to go through them and eliminate ideas that do not work out in practise. The engineer can use their time on decision making instead of generating solutions. Of course the engineer can also add their own ideas and solutions if they notice LLM has missed something. \textcite{amasantiImpactAIGeneratedSolutions2026} argue that software engineers are needed to make final synthesis










